\section{Experiments}
\label{sec:experiments}

\subsection{Setup}
BraTS 2024 glioma~\cite{menze2014brats,baid2021brats}: four co-registered MRI modalities, labels
remapped to $\{$0:Background, 1:NCR, 2:Edema, 3:ET$\}$, patient-wise $75/12.5/12.5$ split
($n_{\text{test}}{=}168$), $z$-scored per modality, $128^3$ inputs, $80$ epochs of AdamW
($10^{-4}$, cosine), mixed precision. The bottleneck is $192\times16^3$. We report a SegResNet
backbone throughout for comparability across configurations; a Swin transformer with a UNETR
decoder reproduces every qualitative finding. Three configurations recur:

\begin{itemize}\itemsep1pt
\item[] \textbf{V0} --- the natural design (residual expert, skips intact).
\item[] \textbf{V2} --- subtractive repair: gated expert, $\alpha{=}0.5$.
\item[] \textbf{DIRECT} --- additive routing, Eq.~\ref{eq:direct}, $\lambda{=}1$.
\end{itemize}

\subsection{The natural design aligns}
\label{sec:aligns}
Table~\ref{tab:aligns} reports the standard picture. Segmentation is competitive, routing tracks
the concepts closely, and on concept-discriminability V0 reaches $0.88$ against $0.74$ for the
strongest post-hoc method, with integrated gradients and Grad-CAM near chance. By every measure
this literature reports, the module works.

\begin{table}[t]\centering\footnotesize
\setlength{\tabcolsep}{3.5pt}
\caption{The natural design by standard metrics (BraTS test, $n{=}168$). Concept-discrim.\ is
concept $k$ vs.\ the \emph{other} tumor voxels; $0.5$ is chance.}
\label{tab:aligns}
\begin{tabular}{lccc}
\toprule
& Dice & $\casfg$ & Concept- \\
& \scriptsize WT/TC/ET & \scriptsize N/E/T & discrim. \\
\midrule
V0 & $0.93/0.88/0.87$ & $0.80/0.96/0.95$ & $\mathbf{0.883}$ \\
\midrule
Occlusion   & --- & --- & $0.741$ \\
Gradient    & --- & --- & $0.601$ \\
Int.\ grad. & --- & --- & $0.583$ \\
Grad-CAM    & --- & --- & $0.521$ \\
\bottomrule
\end{tabular}
\end{table}

Removing $\mathcal{L}_{\text{align}}$ collapses $\casfg$ to $-0.01/0.32/0.09$ while Dice barely
moves ($0.905$ vs.\ $0.927$ WT). We read this at the time as evidence that alignment is
\emph{learned} rather than incidental. It is---but the same number says something we missed: the
segmentation does not need the alignment.

\subsection{\ldots and is inert}
\label{sec:inert}

\paragraph{Intervention.}
Re-routing moves almost nothing. Averaged over concept pairs, $1.1\%$ of re-routed voxels change
label (Table~\ref{tab:repair}, V2 row; V0 gives $0.1\%$ on a smaller pilot). The identity diagonal
is exactly $0$, so this is a property of the model rather than of the instrumentation.

\paragraph{The experts are not degenerate.}
The obvious explanation is false. Experts displace tokens by $0.59$--$0.74$ of their norm and
differ pairwise by $0.93$ (relative $L_2$). The routing selects genuinely different
transformations; they simply do not reach the output.

\paragraph{Bottleneck ablation.}
They do not reach the output because the output does not need them. Replacing the module with raw
features or with \emph{zeros} costs almost nothing: edema Dice $0.879\!\to\!0.865$, ET
$0.862\!\to\!0.806$. A decoder that can segment without its bottleneck will ignore any concept
structure placed there.

\subsection{The probe control}
\label{sec:probe}
A linear probe on the features entering the module scores $0.839$ against the router's $0.872$.
Run on a model trained with \emph{no} concept supervision, the probe still reaches $\mathbf{0.811}$
---above every post-hoc method in Table~\ref{tab:aligns}. The concept structure is largely a
property of a competent segmentation encoder, not of the concept machinery. This has a consequence
beyond our system: reporting that a concept module separates concepts better than saliency
baselines is weak evidence, because logistic regression on an unsupervised bottleneck clears that
bar.

\subsection{The output is a stronger concept classifier than the routing}
\label{sec:segpost}
A second control tightens the same point. The network's own segmentation posterior---free from
the same forward pass---separates the concepts at $0.979$, well above the routing's $0.883$. It
reaches $0.982$ on the transformer backbone despite different pretraining and a different Dice,
so the metric saturates for any competent segmenter: under this window it is close to a
restatement of segmentation accuracy, because ranking voxels is far easier than thresholding
them. Matched to the router's $16^3$ decision budget by average-pooling and re-upsampling through
the identical operator, the posterior falls to $0.935$, so roughly half the residual gap is
spatial resolution rather than concept information.

We report this because omitting it would overstate our case, and because it changes what a
discriminability number means. The posterior is not an explanation: it \emph{is} the prediction,
and restating the output explains nothing about how the output arose. But together with the probe
control (Sec.~\ref{sec:probe}) it establishes that concept-discriminability cannot carry an
interpretability claim on its own---an unsupervised linear probe clears the post-hoc baselines,
and the model's own output beats everything.

\subsection{Repair}
\label{sec:repair}

\begin{table*}[t]\centering\small
\caption{Causal control and its price (BraTS test, $n{=}168$). Label-flip is the mean off-diagonal
fraction of re-routed voxels whose predicted label became the target concept, tumor-restricted and
scope-matched; the identity diagonal is exactly $0$ in every configuration. Subtractive repair
(V2) does not work; additive routing (DIRECT) does, for ${\approx}1$ Dice point.}
\label{tab:repair}
\begin{tabular}{llccccc}
\toprule
& Repair & Label-flip $\uparrow$ & 95\% CI & Dice WT/TC/ET & $\casfg$ N/E/T & routing ECE $\downarrow$ \\
\midrule
V0 & none (natural design) & \multicolumn{2}{c}{$0.001$ \; \footnotesize(pilot, $n{=}5$)} & $0.927/0.883/0.869$ & $0.80/0.96/0.95$ & $0.841$ \\
V2 & gated expert, $\alpha{=}0.5$ & $0.011$ & $[0.009,0.014]$ & $0.925/0.872/0.854$ & $0.79/0.96/0.95$ & $0.673$ \\
\textbf{DIRECT} & additive routing & $\mathbf{0.636}$ & $[0.621,0.651]$ & $0.927/0.873/0.858$ & $\mathbf{0.81}/0.96/0.95$ & $\mathbf{0.681}$ \\
\bottomrule
\end{tabular}
\end{table*}

\paragraph{Subtractive repair fails, instructively.}
Attenuating the skips ($\alpha{=}0.5$) leaves every metric within noise of V0: Dice $-0.002$ WT,
$\casfg$ $-0.001$ on edema, and causal control still at $0.011$. A scalar gate is compensable
during training---normalization affine parameters can simply rescale to undo it---so at
convergence it constrains nothing.

Removing the expert residual is worse than ineffective: pairwise expert divergence \emph{collapses}
from $0.93$ to $0.16$. With $E_k(h)=h+c_k(h)$ each expert supplies only a concept-specific
correction and is free to differ; with $E_k(h)=c_k(h)$ every expert must independently reconstruct
whatever the decoder needs, so they converge on approximately the same map. Making the expert the
sole path made the experts interchangeable.

\paragraph{Additive routing works.}
Table~\ref{tab:repair} gives the result. Causal control rises from $0.011$ $[0.009,0.014]$ to
$\mathbf{0.636}$ $[0.621,0.651]$, a paired difference of $+0.625$ $[+0.609,+0.641]$ with Wilcoxon
$p{=}2.6\times10^{-29}$ over $168$ matched cases. Fig.~\ref{fig:intervention} shows the same result
as images.

Two independent measurements corroborate it. Removing only the routing term collapses the model
(Table~\ref{tab:necessity}): NCR $-99\%$, ET $-80\%$, edema $-53\%$, against ${\approx}1.5\%$ for
zeroing V0's entire bottleneck. And the magnitude ratio $\rho$ converged at $0.35$, so the routing
term dominates the logits roughly $3\!:\!1$---the failure mode we instrumented for did not occur.

\begin{table}[t]\centering\small
\caption{Bottleneck necessity for DIRECT (Dice, $n{=}40$). Each path is ablated separately;
\emph{zero\_routing} is the direct-mode number. Ablating only the decoder's bottleneck input would
leave the routing term intact and wrongly report the bottleneck as unnecessary.}
\label{tab:necessity}
\begin{tabular}{lcccc}
\toprule
& base & bypass & zero\_expert & \textbf{zero\_routing} \\
\midrule
NCR   & $0.483$ & $0.484$ & $0.053$ & $\mathbf{0.006}$ \\
Edema & $0.898$ & $0.477$ & $0.333$ & $\mathbf{0.418}$ \\
ET    & $0.852$ & $0.597$ & $0.428$ & $\mathbf{0.172}$ \\
\bottomrule
\end{tabular}
\end{table}

\paragraph{The price.}
Whole tumor is free ($0.9273$ vs.\ $0.9272$); TC and ET cost ${\approx}1$ Dice point each.
Alignment is marginally \emph{better} than the inert baseline ($\casfg$ NCR $0.809$ vs.\ $0.802$),
calibration improves substantially (ECE $0.841\!\to\!0.681$), and token-level AUROC rises on edema
($0.638\!\to\!0.900$) and ET ($0.622\!\to\!0.839$) while falling on NCR ($0.937\!\to\!0.783$).

\paragraph{Alignment emerges from the task.}
One consequence deserves emphasis. Because the routing logits are part of the output,
$\mathcal{L}_{\text{seg}}$ itself trains the routing to be concept-aligned. By epoch 10, with
$\mathcal{L}_{\text{align}}$ still switched off, $\casfg$ has reached $0.41/0.42$ on edema and ET,
where V0 remains at the floor ($0.08/0.12$). This directly answers the objection that $\casfg$ only
measures what $\mathcal{L}_{\text{align}}$ trains: here alignment appears without it.

\begin{figure}[t]\centering
\includegraphics[width=\linewidth]{figures/fig_intervention.png}
\caption{\textbf{Causal intervention, rendered.} Top: input, ground truth, routing argmax,
baseline prediction. Grid: rows are the \emph{source} concept, columns the \emph{target}; each cell
is the segmentation after forcing every token the router assigned to the row concept through the
column concept's expert, with the re-routed region outlined in white. The diagonal is the identity
control and is pixel-identical to the baseline (changed-voxel fraction $0$ for all three), so the
off-diagonal changes are attributable to re-routing alone.}
\label{fig:intervention}
\end{figure}

\begin{figure}[t]\centering
\includegraphics[width=\linewidth]{figures/fig_intervention_matrices.png}
\caption{\textbf{Causal control, V2 vs.\ DIRECT}, on a shared colour scale. Rows are the source
concept, columns the target; each cell is the fraction of re-routed voxels whose predicted label
became the target. Subtractive repair leaves the matrix uniformly dark; additive routing does not.
Diagonals are the identity control and are exactly zero in both.}
\label{fig:matrices}
\end{figure}

\paragraph{The effect is specific and local.}
A large flip rate would mean little if the intervention simply perturbed everything, so we report
the two controls that separate concept-specific control from generic disturbance
(Table~\ref{tab:specificity}). Classes other than the source and target move by $0.079$ on
average---an order of magnitude below the $0.636$ target effect---and voxels \emph{outside} the
re-routed region move by $0.005$. The source class falls by $0.544$, mirroring the target's rise.
The intervention does what it is supposed to do, where it is supposed to do it.

\begin{table}[t]\centering\footnotesize
\setlength{\tabcolsep}{4pt}
\caption{Specificity and locality of the intervention (DIRECT, $n{=}168$, mean over off-diagonal
concept pairs). Restricting the intervention to the tumor neighborhood does not merely lower the
headline---it makes the effect \emph{cleaner} on both controls.}
\label{tab:specificity}
\begin{tabular}{lcc}
\toprule
& tumor-restricted & all tokens \\
\midrule
target class $\uparrow$              & $\mathbf{0.540}$ & $0.798$ \\
source class $\downarrow$            & $-0.544$ & $-0.514$ \\
off-target classes $\to 0$           & $\mathbf{0.079}$ & $0.151$ \\
outside the region $\to 0$           & $\mathbf{0.005}$ & $0.016$ \\
\bottomrule
\end{tabular}
\end{table}

\paragraph{Scope matters, and the smaller number is the honest one.}
A router typically assigns most of the volume to background: here roughly a quarter of all tokens
go to the edema expert, the majority of them normal tissue. Re-routing \emph{every} token of a
concept therefore repaints the whole brain, and reports $0.816$ rather than $0.636$. That larger
figure is not wrong, but it mostly demonstrates that the output follows the routing everywhere,
which under Eq.~\ref{eq:direct} is close to tautological. Restricting to the tumor neighborhood
measures concept control where it is clinically meaningful---and, as Table~\ref{tab:specificity}
shows, roughly halves the off-target disturbance and cuts the out-of-region leakage by
$3.5\times$. We quote the restricted number throughout and report the scope with every
measurement.

\paragraph{Where it is weakest, and why.}
The effect is not uniform, and the asymmetry is mechanistically informative rather than merely a
caveat. Necrotic core is weak as a \emph{source} ($0.28$--$0.31$) but works as a \emph{target}
($0.51$--$0.66$): the model can be pushed into calling a region necrosis, but pushing necrosis
tokens elsewhere works less well. The source-class posteriors say why. Routing edema or ET tokens
away drops their own class by $0.81$ and $0.58$ respectively, whereas routing necrosis away drops
it by only $0.23$---because the model was never confident about necrosis to begin with (its
isolated Dice is $0.483$, against $0.898$ for edema). There is little confidence to remove, so the
re-routed region is contested rather than reassigned, and off-target classes absorb the difference
($0.09$--$0.11$ from necrosis against $0.03$--$0.06$ elsewhere). Causal control is therefore
bounded by how decisively the model held the original concept, which is a property of the
underlying segmentation problem rather than of the routing mechanism. Necrotic core is the
scarcest concept at a $16^3$ bottleneck---tens of tokens per case---and is the consistent weak
spot across every metric we report, not only this one.

\subsection{The protocol, applied}
\label{sec:summary}
Table~\ref{tab:protocol} collects the three tests across the three configurations. Read down a
column and the story of each design is visible in three numbers; read across a row and the tests
disagree with the alignment metrics in exactly the way Sec.~\ref{sec:protocol} predicts. V0 and V2
are indistinguishable on every causal test while differing on none of the alignment ones, which is
the whole point: alignment cannot separate them, and intervention can.

\begin{table}[t]\centering\footnotesize
\setlength{\tabcolsep}{4pt}
\caption{The diagnostic protocol across configurations. Alignment separates none of them; the
causal tests separate DIRECT from both. Ablation reports edema Dice with the module's contribution
removed, relative to that configuration's own baseline.}
\label{tab:protocol}
\begin{tabular}{lccc}
\toprule
& V0 & V2 & \textbf{DIRECT} \\
\midrule
\multicolumn{4}{l}{\emph{alignment --- cannot separate them}}\\
$\casfg$ (edema)                & $0.962$ & $0.961$ & $0.962$ \\
concept-discrim.                & $0.883$ & $0.879$ & $0.881$ \\
\midrule
\multicolumn{4}{l}{\emph{causal --- separates DIRECT}}\\
T1 intervention $\uparrow$      & $0.001$ & $0.011$ & $\mathbf{0.636}$ \\
T2 ablation, edema $\downarrow$ & $-1.6\%$ & $-0.7\%$ & $\mathbf{-53\%}$ \\
T3 probe vs.\ router            & \multicolumn{3}{c}{$0.839$ vs.\ $0.872$; $0.811$ unsup.} \\
\bottomrule
\end{tabular}
\end{table}

\subsection{Backbone transfer}
\label{sec:backbone2}
Table~\ref{tab:backbone} reports the natural design on both backbones. A Swin transformer with a
UNETR decoder and a SegResNet CNN differ in pretraining, capacity and accuracy, yet produce the
same qualitative picture: strong alignment, concept-discriminability far above the post-hoc
baselines, and the same inertness. The leaks are properties of the architecture family rather than
of one backbone, and the transformer exhibits the placement failure of Sec.~\ref{sec:leaks} in its
sharpest form---its module sits at stage~2 of~5, so the two deeper and more semantic encoder
stages bypass it entirely.

\begin{table}[t]\centering\footnotesize
\setlength{\tabcolsep}{4pt}
\caption{The natural design on two backbones. Both align well; both are inert.}
\label{tab:backbone}
\begin{tabular}{lccc}
\toprule
& Dice WT/TC/ET & $\casfg$ N/E/T & discrim. \\
\midrule
Swin + UNETR & $0.915/0.845/0.837$ & $0.69/0.94/0.93$ & $0.872$ \\
SegResNet    & $0.927/0.883/0.869$ & $0.80/0.96/0.95$ & $0.883$ \\
\bottomrule
\end{tabular}
\end{table}
